% Version 2022-09-20
% ISPRS full-paper template (isprs.cls) -- content: Levin Giersch project report
% Requires isprs.cls and the bibliography file Bib.bib in the same folder.
% Compile with: pdflatex -> bibtex -> pdflatex -> pdflatex

\documentclass{isprs} % isprs class modified 23-04-2019 (Dennis Wittich)
\usepackage{subfigure}
\usepackage{setspace}
\usepackage{geometry} % added 27-02-2014 Markus Englich
\usepackage{epstopdf}
\usepackage[labelsep=period]{caption}  % added 14-04-2016 Markus Englich - Recommendation by Sebastian Brocks
\usepackage[british]{babel}
\usepackage[hang]{footmisc}
\def\footnotemargin{1em} % added 08-01-2020 Dennis Wittich

% --- added so the special characters in the report content render ---
\usepackage{textcomp} % provides \textdegree
\usepackage{amsmath}
\usepackage{booktabs} % \toprule, \midrule, \bottomrule for the dataset table
\usepackage{url} % defines \url, used by isprs.bst for web/DOI references
% --------------------------------------------------------------------

\geometry{a4paper, top=25mm, left=20mm, right=20mm, bottom=25mm, headsep=10mm, footskip=12mm} % added 27-02-2014 Markus Englich

\captionsetup{justification=centering,font=normal} % thanks to Niclas Borlin 05-05-2016
\captionsetup[figure]{font=small} % added 23-04-2019 Dennis Wittich
\captionsetup[table]{font=small} % added 23-04-2019 Dennis Wittich

\begin{document}

\title{A harmonised point-extraction and comparison tool for heterogeneous paleoclimate datasets}
\date{}

\author{Levin Giersch}

\address{
	Faculty VI - Planning Building Environment, Geodesy and Geoinformation Science, M.Sc. - l.giersch@campus.tu-berlin.de\\
	Student ID: 1690 0515040, Date: \today, Lecturer: Prof. Dr.-Ing. Martin Kada
}

\abstract{

- paleoclimate simulations are very important for understanding earth systems
- but they are distributed in diverse and sometimes incompatible formats
	- different time orgigin and unit
	- spatial grid
	- conventions
	- temperature variables
	- time and spatial resolutions

- this project provides a reproducable python toolchain to download preprocess, and compare temperatre archives for widely used datasets:
	- TraCE-21k,
	- CHELSA-TraCE21k-centennial,
	- Beyer2020/2021
	- PalMod2
- harmonise to kiloyears-before-present (ka BP) scale
- harmonise spatial grids
- convert temp units to degree celsius
- extract temperature at user-defined point (use berlin as example)

- globally extract temperature of each dataset at each grid cell size of trace21k (most coarse), compare pair-wise

- all models agree on cold period around 20 kaBP followed by
- but disagree in absolute temperature by several degrees.

- tried to find pattern in differences between dataset pair: global maps plotting the temp difference between each dataset-pair:
	- largest differences in high norther latitudes

- this report documents datasets, data aquisition method, preprocessing, harmonisation, data extraction at point and global comparison

- project is organized in several jupyter notebooks creating a full automatic pipeline \cite{giersch2026_paleoclimate_model_comparison}
}

\keywords{Paleoclimate, Data harmonisation, TraCE-21k, CHELSA, PalMod2, Temperature comparison}

\maketitle

\pagestyle{plain} % page number centered at the bottom of every page

\sloppy

%-------------------------------------------------------------

\section{Introduction}
\subsection{Why comparing paleoclimate models matters}

- Paleoclimate models are useful for:\\
	- Reconstructing Earth's climate before direct measurements existed\\
	- Understanding how the climate responds to forcings like solar output, volcanic eruptions, orbital cycles, and CO2\\
	- Testing and validating modern climate models against known past states\\
	- Establishing the natural variability range of climate, separate from human influence\\
	- Constraining climate sensitivity (temperature change per CO2 doubling)\\
	- Studying planetary climate evolution -- Mars, Venus, and other terrestrial bodies\\
	- Using past warm periods (e.g. Pliocene) as analogues for near-future climate projections\\

- there are many different paleoclimate models.\\
- coparing them is important to detect outliers, irregularities or strong differences eg. to help improve the models\\
- it is interesting to see where they disagree by how much
- systematic differences show where the models are robust and where they diverge

\subsection{Aim of this project}
- harmonise diverse and incompatible paleoclimate datasets
- tool to extract temperature (air and soil) values at user difined point
- plot temp values over timeseries
- compare paleoclimate models pairwise globaly at resoulution of trace-21k grid size and plot results as map (mean temp differences)



\section{The Data}
- four datasets

\begin{table*}[ht]
\centering
\footnotesize
\begin{tabular}{llllll}
\toprule
Dataset & Variables & Time coverage & Time step & Spatial resolution & Format \\
\midrule
Beyer2020 (v4)              & air temp            & 120--0 ka BP & 1--2 kyr & 0.5\textdegree{}                       & NetCDF \\
PalMod2                     & air temp, soil temp & 25--0 ka BP  & annual   & \textasciitilde3.75\textdegree{} (T31) & NetCDF \\
TraCE-21k                   & air temp, soil temp & 22--0 ka BP  & decadal  & \textasciitilde3.75\textdegree{} (T31) & NetCDF \\
CHELSA-TraCE21k-centennial  & air temp            & 21--0 ka BP  & 100 yr   & \textasciitilde1 km (0.01\textdegree{}) & GeoTIFF $\rightarrow$ NetCDF \\
\bottomrule
\end{tabular}
\caption{Overview of the four paleoclimate datasets.}
\label{tab:datasets}
\end{table*}

\subsection{Trace21k}
- Simulation of the Transient Climate of the Last 21,000 Years
- Chosen Version: TraCE Main Land Decadal Average\\
- developed by: 2021 \\
- year of development: see source\\
- how where they made: fully-coupled, non-accelerated atmosphere-ocean-sea ice-land surface CCSM3 simulation\\
- what was the goal: \\
- nr. files: 2\\
- size (gb): 449,2 MB\\
- type: netcdf\\
- used variables: TSA, TSOI\\
- how to get: (via python script in my data\_downloader.ipynb notebook)\\
- other info:  \\
- extend time: 22,000 BCE to 1990 CE \\
- resolution time: decadal averages\\
- extend space: global land\\
- resolution space: T31\_gx3\\
- source: \cite{trace21k_data}\\

\subsection{CHELSA-TraCE21k}
- Chosen Version: CHELSA-TraCE21k-centennial\\
- developed by: see source\\
- year of development: see source \\
- how where they made: generated with the CHELSA-TraCE21k downscaling model. \\
- what was the goal: \\
- nr. files: \\
- size (gb): \\
- type: geotiff -> need to be converted to netcdf\\
- used variables: tasmin, tasmax -> tasmin+tasmax/2 to get mean\\
- how to get: (via python script in my data\_downloader.ipynb notebook)\\
- other info: \\
- extend time: Start Time 21k BP, End Time 0 BP\\
- resolution time: monthly averages, 100 year timesteps\\
- extend space: global \\
- resolution space: kilometer scale\\
- source: \cite{chelsa_trace21k_paper}, \cite{chelsa_trace21k_data}, \cite{chelsa_trace21k_paper}\\


\subsection{Beyer 2021 (v4)}
- Chosen Version: 12293345 version 4 (figshare, published 24 June 2021)\\
- developed by: Robert M. Beyer, Mario Krapp \& Andrea Manica (Department of Zoology, University of Cambridge)\\
- year of development: original paper published 14 July 2020 (received 03 July 2019); dataset v4 published 24 June 2021; addendum paper published 29 September 2021\\
- how were they made: Delta-method-based downscaling and bias correction, combining medium-resolution HadCM3 transient simulations with high-resolution HadAM3H snapshots and present-day observational data\\
- what was the goal: Fill the gap between high-temporal-resolution but low-spatial-resolution palaeoclimate simulations and high-spatial-resolution snapshot-only datasets\\
- nr. files: \\
- size (GB): \\
- type: NetCDF (data) + R / Python / MATLAB scripts\\
- used variables: \\
- how to get: (via python script in my data\_downloader.ipynb notebook)\\
- other info: v4 adds global monthly NPP relative to v3; documented in the 2021 addendum\\
- extend time: 120000 BP to pre-industrial modern era (0 BP)\\
- resolution time: 2000-yr steps from 120000--22000 BP; 1000-yr steps from 22000 BP to pre-industrial (72 snapshots total)\\
- extend space: global terrestrial (full 360\textdegree{} longitude; latitude coverage 720x 300 cells at 0.5\textdegree{}, i.e. 150\textdegree{} of latitude)\\
- resolution space: 0.5\textdegree{} equirectangular grid (\textasciitilde55 km at the equator)\\
- source: \cite{beyer2020_paper}, \cite{beyer2021_addendum}, \cite{beyer2021_data_4}, \cite{beyer2020_data_3}\\


\subsection{PalMod 2}
- Chosen Version: PalMod2 MPI-M MPI-ESM1-2-CR Transient Simulations of the Last Deglaciation with prescribed ice sheets from GLAC-1D reconstructions (r1i1p1f1)\\
- developed by: Mikolajewicz, U.; Kapsch, M.-L.; Gayler, V.; Meccia, V. L.; Riddick, T.; Ziemen, F. A.; Schannwell, C. (Max Planck Institute for Meteorology / DKRZ)\\
- year of development: dataset DOI published 2023\\
- how where they made: fully-coupled transient simulation with the Max Planck Institute Earth System Model version 1.2 in coarse resolution (MPI-ESM1-2-CR), combining the ECHAM6.3 spectral atmosphere (T31, 31 vertical levels)
- what was the goal: produce a transient, fully-coupled simulation of the Last Deglaciation with realistic, time-evolving ice-sheet, topography and land-sea-mask forcing as part of the PalMod2 project\\
- nr. files: \\
- size (gb): \\
- type: NetCDF\\
- used variables: tas (PMMXMCRTDGr111Amtasgn30201, near-surface air temperature, monthly)
tsl: (PMMXMCRTDGr111Lmtslgn30201, soil temperature by layer, monthly) \\
- how to get: WDCC archive at DKRZ; download via jblob client after registration \cite{wdcc_jblob} (via python script in my data\_downloader.ipynb notebook)\\
- other info: \\
- extend time: 25000 BP to 0 BP (= 1950 CE)\\
- resolution time: variable-dependent --- monthly for the standard atmosphere/land/ocean/sea-ice tables (Amon, Lmon, Omon, SImon, LImon, Emon), decadal (dec, Odec) and fixed (fx, Ofx); GLAC-1D ice-sheet/topography forcing updated every 10 years\\
- extend space: global\\
- resolution space: atmosphere/land T31 (\textasciitilde3,75\textdegree{}), ocean MPIOM1.6 nominal \textasciitilde3\textdegree{}\\
- source: \cite{palmod2_data}\\


%-------------------------------------------------------------
\section{Pre-processing}
- trace-21k and beyer2020 are already analysis ready netcdfs, the do not have to be preprocessed


\subsection{palmod2}
- tas and tsl are delivered as monthly netcdf files
- a annual mean has to be computed from the monthly values (cdo yearmean)
- then merge all the files to a single netcdf file along time axis.
- merged outputs have 25000 annual time steps


\subsection{Chelsa tarce 21k}
- data is provided as montly geotiffs
- use regex on filenames to group files by year
- stack and average the yearly tiffs
- convert from kelvin * 10 to kelvin
- concatenate along time axis
- add metadata
- write as gzip compressed netcdf files per variable (tasmin and tasmax)
- compute tasmean from tasmin and tasmax
- this process costs alot of time and compute power and disk space. ca 1tb disk space, several days computation on modern laptop. cdo usage not possible because of ram limitation so dask has to be used (lazy) \cite{giersch2026_paleoclimate_model_comparison}


%-------------------------------------------------------------
\section{Homogenization}
\subsection{Variables}
- to compare the datasets thy have to share a common variable, time format, space and unit
- homogenisation made in comparison notebook \cite{giersch2026_paleoclimate_model_comparison}

- DATASET.cvs HIER REIN PACKEN ALS TABELLE

\subsection{Time}
- time extend \\
- time resolution \\

\subsection{Space}
- grid sizes \\
- changes in north south \\
- different resolutions for different times \\

%-------------------------------------------------------------
\section{Point Extraction}
- do i just take the grid cell the point is in? \\

%-------------------------------------------------------------
\section{Results)}
- For one example Point \\
- differences between models \\

%-------------------------------------------------------------
\section{Future}
- Auto iterate over points and compare \\
- correlate with eg. north south, climate zones, water vs land etc. \\


%-------------------------------------------------------------
\section{Conclusion}
- hard to compare due to differences eg. resolutions \\
- comparison can helt to find out where the differ the most -> conlusion to how to improve models? \\


%-------------------------------------------------------------
{
	\begin{spacing}{1.17}
		\normalsize
		\bibliography{Bib} % your bibliography file Bib.bib; style is given in isprs.cls
	\end{spacing}
}

\end{document}